% Development methods supplement. No external novelty claim is established.
\documentclass[11pt,a4paper]{article}
\usepackage[T1]{fontenc}
\usepackage{amsmath,amssymb,booktabs}
\usepackage[margin=25mm]{geometry}
\title{Hierarchical discrete adjoints in PhotonWeave\\\large Development methods supplement}
\author{Hyoseok Park}
\date{20 September 2026}
\begin{document}
\maketitle

\section{Scope of the implemented operation}
The current differentiable operation advances nondispersive Yee fields,
including fixed complex Bloch fields in the resident execution path,
and convolutional perfectly matched layer memories on a fixed mesh. The
permittivity can be scalar or diagonal. Point observations and fixed collocated
spectral planes define the current output contract. A geometry map implemented with Torch operations can precede
this operation. This supplement does not establish derivatives of every
forward-solver feature, convergence of physical shape derivatives, or a new
checkpointing algorithm relative to prior work.

Let $s_t$ contain all six field components and the CPML auxiliary states.
The implemented full step updates the electric field and then the magnetic
field, including prepared sources and their prescribed times. Write this as
\begin{equation}
 s_{t+1}=F_t(s_t,\epsilon),\qquad y_t=P_t s_{t+1}.
\end{equation}
For a real objective $L(y_0,\ldots,y_{T-1})$, let $\bar y_t$ denote its incoming
observation derivative. The discrete reverse recurrence is
\begin{align}
 v_{t+1} &= \bar s_{t+1}+P_t^{\mathsf H}\bar y_t,\\
 \bar s_t &= (D_s F_t)^{\mathsf H}v_{t+1},\\
 \bar\epsilon &\mathrel{+}=\operatorname{Re}\left[(D_\epsilon F_t)^{\mathsf H}v_{t+1}\right].
\end{align}
The CUDA transpose includes CPML state derivatives. Reversing lossless Maxwell
fields alone would not reproduce this discrete recurrence in an absorbing
domain. Checkpoint replay reconstructs the complete primal state required by
the transpose. A Torch geometry map $\epsilon=g(\theta)$ then supplies
$\bar\theta=(D_\theta g)^{\mathsf T}\bar\epsilon$.

For fixed Bloch factor $p=\exp(\mathrm{i}\phi)$, the forward seam difference
is $pF_0-F_{N-1}$ and the backward seam difference is
$F_0-F_{N-1}/p$. Their adjoints multiply by conjugate seam coefficients.
Complex CPML memories participate in checkpoint replay. The complex resident
path uses Torch on CPU and optionally fused forward and transpose kernels on
CUDA. Complex spatial slab streaming remains unsupported.
A 20-degree TE slab check at 1.55 micrometres gave transmission error 0.004442
against the Fresnel slab formula and conservation error $9.58\times10^{-6}$.
Its fixed-geometry index derivative agreed with a central difference to relative
error $1.18\times10^{-8}$. This is a single-frequency validation rather than a
full pupil, broadband constant-angle or shape-convergence result.

\section{Bounded spectral observations}
For a fixed frequency $f_j$ and fixed real window $w_t$, a point component has
the time-integral observation
\begin{equation}
 z_{jm}=\Delta t\sum_{t=0}^{T-1}w_t y_{tm}
 \exp\left[-2\pi\mathrm{i}f_j(t+1+\delta_m)\Delta t\right],
\end{equation}
where $\delta_m=0$ for electric components and $\delta_m=1/2$ for magnetic
components under the implemented update ordering. For the incoming complex
derivative $\bar z_{jm}$, the real observation seed is
\begin{equation}
 \bar y_{tm}=\Delta t w_t\operatorname{Re}\sum_j
 \bar z_{jm}\exp\left[2\pi\mathrm{i}f_j(t+1+\delta_m)\Delta t\right].
\end{equation}
The complex derivative convention satisfies
$dL=\operatorname{Re}\sum_{jm}\bar z_{jm}^{*}dz_{jm}$.
For complex time samples, the observation seed omits the real-part projection.
Both the accumulation and its transpose operate on bounded time blocks.
The point-signal history and a full time-by-frequency matrix are not retained.
With $F$ frequencies, $M$ observations and block length $B$, the additional
workspace scales as $O(FM+FB+BM)$. A supplied window and the prepared sources
still require time-length arrays. These observations are point fields, not
collocated plane flux or normalized port transmission. They establish neither
a new Fourier algorithm nor completion of the color-router objective.

\section{Collocated spectral planes}
For each component, fixed trilinear interpolation maps indexed Yee samples to
physical quadrature points. Shared support samples are deduplicated across
planes. Interpolation commutes with the linear Fourier accumulation because
its weights are time independent. The plane API uses the positive exponential
of the native frequency monitor, obtained by conjugating the point transform
kernel. Its transpose uses the same conjugated convention.

For a plane with coordinate normal $a$, let $(a,b,c)$ be a cyclic ordering.
The integrated reduced spectral power is
\begin{equation}
 P_a=\frac{1}{2}\sum_q A_q\operatorname{Re}
 \left(E_b(q)H_c(q)^*-E_c(q)H_b(q)^*\right),
\end{equation}
where $A_q$ is the clipped physical-cell quadrature weight. Torch transposes
the interpolation and power objective into spectral sample derivatives.
The field solver then applies its bounded discrete adjoint. Matched-reference
normalization divides by absolute incident flux. Reflection subtraction acts
on the complex fields before power is formed. This is not mode decomposition.

A fixed dielectric slab with index 1.5 and thickness 0.2 micrometres was tested
at nine wavelengths between 1.3 and 1.8 micrometres. At a 0.025 micrometre mesh
and 1,600 steps, the maximum Fresnel transmission error was 0.001535 and the
maximum error in $T+R=1$ was $4.44\times10^{-6}$. The mean-transmission index
derivative agreed with a central difference of step $10^{-4}$ to relative
error $1.14\times10^{-8}$. This validates one fixed-mask material derivative,
not sharp-interface shape convergence or the color-router objective.

\section{Joint Gaussian target objectives}
The real local measurement model is $Y=AX+N$. Explicit joint moments
$\Sigma_{XX}$, $\Sigma_{XZ}$ and $\Sigma_{ZZ}$ retain target uncertainty
outside the sampled scene latent. Let $L_XL_X^{\mathsf T}=\Sigma_{XX}$ and
$L_NL_N^{\mathsf T}=\Sigma_N$. Define
\begin{align}
 C&=L_X^{-1}\Sigma_{XZ},& W&=L_N^{-1}AL_X,\\
 R&=\Sigma_{ZZ}-C^{\mathsf T}C,& Q&=I+W^{\mathsf T}W.
\end{align}
The implemented posterior and information are
\begin{align}
 \Sigma_{Z|Y}&=R+C^{\mathsf T}Q^{-1}C,\\
 I(Y;Z)&=\frac{\log\det\Sigma_{ZZ}-\log\det\Sigma_{Z|Y}}{2\log 2}.
\end{align}
Triangular and Cholesky solves replace explicit inverses. A response-dependent
shot covariance remains differentiable. Invalid covariance inputs are rejected
without jitter. This standard Gaussian model is not a novelty claim or a
replacement for the locked optical response, statistical prior and electron
calibration in the color-router study. Synthetic comparisons verify algebra
and derivatives only.

\section{Slab ownership and the transpose}
A time block contains $K$ full steps. Each x slab retains the full y and z
extent and a halo of $K$ cells on either x side. The electric update uses a
backward x difference with support $\{i-1,i\}$, while the magnetic update
uses a forward difference with support $\{i,i+1\}$. Their composition has
support $\{i-1,i,i+1\}$. Pointwise ADE recurrences do not enlarge this set.
Thus a complete step has x radius one, and $K$ complete steps have radius $K$.
Every slab reads the same immutable
block-start state. Only its owned interior is committed to the next global
state. CPML ownership follows the target cell of its derivative rather than
an unshifted field index. Periodic halos can contain repeated physical cells.

Let $S_i$ gather tile $i$, including halo replication. Let $R_i$ scatter only
its owned output and let $G_i$ be its local time-block operation. The global
block map has the form
\begin{equation}
 s_{b+1}=\sum_i R_i G_i(S_i s_b,S_i^\epsilon\epsilon).
\end{equation}
The input derivative accumulates complete halo contributions,
\begin{equation}
 \bar s_b=\sum_i S_i^{\mathsf T}
 (D_sG_i)^{\mathsf T}R_i^{\mathsf T}\bar s_{b+1},
\end{equation}
with additional observation terms at each local time. Permittivity derivatives
use the corresponding gather transpose $(S_i^\epsilon)^{\mathsf T}$. Repeated
periodic copies are summed. Overwriting halo derivatives would omit valid
dependency paths. Each observation is emitted and seeded only by its owner.

\section{Two levels of bounded replay}
Global block checkpoints reside in host DRAM or immutable file banks. Local checkpoints reside in the
active tile workspace and include E, H and every local CPML state. These are
different budgets. A fixed number of local checkpoints does not retain a
time-length field graph on the GPU.

For a schedule of $n$ operations and $q$ additional checkpoints, denote the
number of forward replay operations by $A(n,q)$. The current recursive split
$m=m(n,q)$ obeys
\begin{align}
 A(1,q)&=0,\\
 A(n,0)&=n(n-1)/2,\\
 A(n,q)&=m+A(n-m,q-1)+A(m,q),\qquad q>0.
\end{align}
This describes the implemented split and is not an optimality theorem. It is
used once for global blocks and again for local steps. A zero-checkpoint tile
requires $K(K-1)/2$ replay steps. Additional checkpoints reduce replay in the
measured deeper-block case but add state copies and allocations. Consequently,
fewer replay steps do not imply a lower complete-iteration time.

For $N$ global cells, maximum extended tile volume $V$, $p$ asynchronous slots,
and fixed checkpoint counts $q_g,q_\ell$, the physical-state storage scales as
\begin{align}
 M_{\rm host}&=O\bigl(N(q_g+1)+pV(q_\ell+1)\bigr),\\
 M_{\rm device}&=O\bigl(pV(q_\ell+1)\bigr).
\end{align}
The constants include state banks, adjoints, coefficients and staging buffers.
The all-zero host initial fields and CPML memories use scalar-backed views.
They are read by tile extraction and are not advanced in place. The host
template omits inverse permittivity because each active tile constructs its
own coefficients. Evolved global banks and their adjoints remain dense in
DRAM. This reduces a persistent initialization constant rather than changing
the asymptotic storage of the evolved states or exploiting sparse fields.
Prepared drives and returned time traces also require time-dependent storage.
Geometry graphs, optimizer allocations, CUDA context and allocator caching
are outside the solver reservation. These equations therefore do not imply
a time-independent bound on the entire process.

\section{Selection and current evidence}
The measured policy selector compares admitted slab sizes, block depths,
transfer policies and optional local checkpoints. For requested probe length
$P$ and candidate depth $K_i$, the first calibration duration is
$L_i=\min(T,K_i\lceil P/K_i\rceil)$ and the second is
$\min(T,2L_i)$. Thus a truncated block occurs only when measuring the complete
target run. Two distinct durations separate startup and block costs. Every
distinct duration is checked against the same admitted reference policy.
Reference storage and additional evaluations count toward memory admission and
calibration cost. An explicit maximum calibration length bounds the search.
An optional refinement remeasures the two initially fastest predicted policies
at one additional doubled duration when the target and calibration limit allow
it. The two longest measurements then determine their revised predictions.
The initial estimates and all additional costs are retained in the report.
This bounded refinement does not establish statistical confidence or global
optimality.
Refinement is disabled by default because sequential development trials
selected the same fastest candidate with and without the additional measurements.
The actual global replay count extrapolates the full iteration. Startup noise
and a partial terminal block in the target can still bias that prediction,
so its quality must be tested at the intended duration.

On an RTX 5880 Ada, the $128\times32\times32$ FP64 case with 64 steps and
$K=32$ required 0.591 seconds per complete streamed iteration without local
checkpoints and 0.396 seconds with one. These are medians of three measurements
after a warm-up. The resident solver required 0.073 seconds in the same test.
A smaller $64\times16\times16$ case with $K=12$ regressed when local checkpoints
were enabled. The measurements support an optional memory--replay trade-off,
not universal acceleration or superiority over another solver.

The executable benchmark, raw repetitions, source hashes, gradient checks and
memory scope are recorded in the companion local replay validation report.
An experimental buffered-file policy stores global primal and adjoint banks
outside application-owned DRAM. Slab reads and writes replace whole-bank
allocation and checkpoints retain immutable file versions. The material and
its gradient remain CPU tensors. The OS page cache is not included in the host
reservation. File cleanup is included in measured iteration time.
For a $256\times256\times128$ FP32 case with 16 steps, one timed iteration
after a warm-up took 3.142 seconds with DRAM banks and 8.210 seconds with file
banks on an RTX 5880 Ada. Gradients agreed exactly. The host reservations were
3.314 and 0.891 GiB, respectively. These are allocation estimates rather than
measured resident memory. This experiment establishes a lossless file-backed
execution path, not cold-storage throughput or acceleration.

Large-domain physical accuracy, sustained storage-backed performance, fully
automatic tier selection and the color-router objective remain separate
validation requirements.

\section{Coupled objectives across independent cases}
For fixed cases $f_i$ and shared design $p$, let
$L=\ell(f_1(p),\ldots,f_m(p))$. The runtime first evaluates the cases
without recording their graphs and retains their compact outputs. Backward
recomputes case $i$ and applies its objective seed $\bar y_i$, giving
\[
\nabla_p L=\sum_{i=1}^{m} (D f_i(p))^{\mathsf T}\bar y_i.
\]
For complex outputs the real differential uses the corresponding Hermitian
vector--Jacobian convention. Only one case graph is retained at a time.
The packed outputs, shared design, accumulated gradient and objective still
require memory. Cases must be deterministic and unchanged between passes.
An output replay check detects value drift but cannot prove purity.

Eight fixed complex Bloch cases on a $64^3$ grid with 24 steps reduced peak
Torch CUDA allocation from 358,044,672 to 159,702,528 bytes on an RTX 5880 Ada.
Median complete iteration time increased from 4.161 to 6.778 seconds over
three alternating repetitions after warm-up. Scalar gradients agreed exactly.
This establishes a memory--recomputation trade-off for a synthetic coupled
objective. It does not establish color-router fidelity or runtime superiority.

\section{Detector allocation proxy}
A fixed midpoint quadrature can be specified independently of the field mesh.
The collocated electric samples define
\[
Q_j=\sum_{q\in j} w_q \sum_{c\in\{x,y,z\}} |E_c(q)|^2,
\qquad a_j=T\frac{Q_j}{\sum_k Q_k}.
\]
Both the electric field and supplied total transmission $T$ differentiate.
These entries allocate transmission and do not represent local absorption or
charge collection. For an incoherent pupil, allocation is computed separately
for each ray and polarization before weighted averaging. The reference uses
power-normalized RCWA transmission. A matched FDTD flux reference requires
independent physical validation before claiming an equivalent CR response.

\section{Fixed spectral polarization synthesis}
Two independent source responses in a homogeneous reference define the
transverse matrix $M(f)$. Its columns are quadrature averages of the complex
tangential electric field after removing the prescribed transverse Bloch phase.
For desired tangential polarization $e_t(f)$, fixed coefficients satisfy
\[
M(f)c(f)=e_t(f),\qquad
E(f)=\sum_{j=1}^{2}c_j(f)E_j(f),\quad
H(f)=\sum_{j=1}^{2}c_j(f)H_j(f).
\]
The same coefficients combine reference and structure fields before power is
computed. Structure derivatives traverse both source solves. The coefficients
are fixed and do not supply source or angle derivatives. The method assumes a
homogeneous forward-wave reference and rejects an ill-conditioned basis.
It is a selected-frequency synthesis rather than an oblique one-way source.

\section{Equivalent-electron objective composition}
For spectral response $R$, mean irradiance $s$, quantum efficiency $q$,
quadrature $w$, and fixed calibration $c$, define
\[
U_{ij}=R_{ij}s_j q_j w_j\frac{\lambda_j}{hc_0}c,
\qquad A=UB^{\mathsf T},\qquad \mu=U\mathbf{1}.
\]
Here $B$ is the explicitly supplied latent spectral basis. Exposure factor
$e$ gives $A_e=eA$ and mean $e\mu$. The Gaussian shot--read approximation
uses $N_e=\operatorname{diag}(e\mu+\sigma_r^2)$, retaining the response
contribution to the noise derivative. Fixed exposure probabilities average
information per raw pixel. The design is not renormalized to a fixed output
brightness. A locked development-context comparison on three generated optical
responses gave maximum score error $6.0\times10^{-15}$ and response-gradient
relative error $4.81\times10^{-14}$ across CPU and CUDA evaluations.
This validates objective composition rather than an optical CR reconstruction.

\section{Periodic density transfer and dense observations}
Piecewise-constant transverse density pixels are transferred by exact box
intersection weights with periodic wrapping. Each electric component receives
an arithmetic volume average at its own Yee location. The transfer preserves
integrated density and is linear in density values. It is not a conformal
interface construction and requires optical mesh convergence.

Dense planes exposed an observation overhead independent of the Yee update.
Prepared per-family indices replace per-sample Python indexing with batched
gather and adjoint scatter-add, retaining duplicate-observation accumulation.
Index storage is included in admission estimates. For 10,000 complex E/H
samples on RTX 3060, the observation-only median decreased from 51.064 to
0.232 milliseconds over five alternating repetitions. This measurement does
not establish an end-to-end solver acceleration factor.

\section{Patterned-layer density derivative pilot}
A hash-identified binary pattern was relaxed to $0.01+0.98\rho$ to permit
central perturbations inside the density interval. A selected-ray test on a
$40\times40\times166$ grid with 1600 FP64 steps combined two calibrated source
bases and normalized detector allocation. Its weighted quadratic allocation
objective is a derivative probe rather than the full CR information objective.
For a smooth transverse density direction, the adjoint derivative was
0.219224692705. Central differences at steps $10^{-3}$ and $5\times10^{-4}$
had relative discrepancies $1.223\times10^{-5}$ and $3.057\times10^{-6}$.
The approximately fourfold reduction is consistent with second-order
truncation. A complete iteration on RTX 3060 took 158.54 seconds and peaked
at 332,274,176 bytes of Torch CUDA allocation with four checkpoint slots
and sequential source-case replay. This validates one discrete directional
derivative and does not establish physical gradient convergence.

\section{Spectral and pupil assembly}
For fixed wavelength $\lambda_\ell$ and pupil ray $r$, the detector response
$R_{\ell r p c}$ retains polarization $p$ and channel $c$ as separate axes.
The incoherent spectral response is
\[
 A_{c\ell}=\sum_r w_r\frac{1}{P}\sum_{p=1}^{P}R_{\ell r p c}.
\]
The fixed intensity weights $w_r$ are not renormalized. Consequently,
illumination falloff and exact-weight pupil subsets are preserved. Compact
case outputs are retained while backward replays one wavelength--ray graph
at a time. The selected-layer implementation recalculates homogeneous
references without gradients on every call and retains two source graphs
within the active case. This is a compute--memory tradeoff rather than a
parallel throughput claim. The assembled response feeds the spectral electron
operator and exposure-weighted information objective. Full CR optical
convergence and optimization remain unverified.

\section{Bounded homogeneous-reference reuse}
An optional CPU cache retains compact homogeneous spectral-plane responses
under a tensor-byte budget. Least-recently-used eviction limits retained
storage. Keys include the fixed simulation, frequency, quadrature, precision
and device. Density changes require a new sample solve but do not invalidate
these homogeneous references. Cached returns are independent copies, and no
full field history is retained. The cache can eliminate repeated reference
solves in backward when entries remain resident. It does not reduce the first
forward sweep and no measured speedup is claimed at this stage.

\section{Complex fused forward execution}
A selectable resident CUDA path represents complex fields by two adjacent
real lanes. Each lane evaluates Yee differences and CPML recurrence in a
single kernel. Seam samples use the complex Bloch factor for the forward
difference and its reciprocal for the backward difference. Real diagonal
material coefficients scale the resulting curl. Separate H and E launches
preserve the global update order, with source injection between them.
The implementation avoids full-volume derivative temporaries during forward
and replay. The adjoint transpose independently selects Torch or the fused
complex implementation below. Tests compare
both precisions, nonuniform metrics, all CPML memories and density VJPs,
including three-dimensional Bloch seams and asynchronous checkpoint replay.
This forward kernel alone does not implement spatial streaming. The complex
tile integration is described below.

For one relaxed CR seed at 540 nm and one pupil ray, the complete FP64 forward
API was measured at 50 nm with 1600 steps on RTX 3060. Three alternating
measurements followed one warmup per backend. The median decreased from
31.4499 seconds to 8.3904 seconds, a factor of 3.7483, with a final response
difference of $5.55\times10^{-17}$. Peak Torch CUDA allocation decreased from
143,787,520 to 88,134,144 bytes. Both paths include two homogeneous references,
two sample solves and detector reduction. These measurements do not include
backward or establish physical convergence or superiority over another solver.

\section{Fused complex transpose}
A separate optional CUDA transpose gathers conjugated Bloch seam contributions
and updates complex CPML adjoints using two buffers. One thread owns each
cell's real material derivative,
\[
 \bar{\epsilon} \mathrel{+}=
 -s\,\operatorname{Re}\!\left(\bar{E}^{*}\operatorname{curl}H\right)
 /\epsilon^2,
\]
where $s$ is the reduced Yee Courant coefficient. The scalar material case sums
the three electric-component contributions. Separate H and E transpose launches
preserve update order. Observation seeds still use prepared Torch gather and
scatter operations. The optional kernel supports first-order real dielectric
VJPs with complex checkpoint replay. Dispersive derivatives and trainable
Bloch phases remain unsupported.

For the selected relaxed CR seed case, matched fused-forward runs with four
checkpoint slots and no reference cache compared Torch and fused transposes.
Three alternating measurements after one warmup per path gave complete
objective/VJP medians of 50.2011 and 32.4018 seconds, respectively. Peak Torch
CUDA allocations were 359,952,384 and 328,486,400 bytes, with a common starting
allocation of 17,170,944 bytes in every measured repetition. The maximum
per-pixel gradient difference was $2.17\times10^{-19}$. These measurements
follow an explicit recursive-closure lifetime fix in resident and tiled
replay. Lifetime tests disable cyclic garbage collection and check object
release after completed result graphs are discarded. This comparison does
not establish full-pupil optimization or a competing-solver advantage.

\section{Complex spatial tiles and measured communication costs}
For an unwrapped tile coordinate $i$ and global extent $N_x$, define the
Bloch extension $P_i=p^{\lfloor i/N_x\rfloor}$, where $p$ is the fixed
unit-modulus phase. The tile reads $P_i F_{i\bmod N_x}$ for electric,
magnetic and transverse CPML states. Repeated source images receive the
same phase. Material coefficients remain periodic without this factor.
The discrete adjoint accumulates $P_i^*\bar F_i$ at the corresponding
global index, including repeated windings. Real material gradients sum
without a phase factor. Global checkpoint replay and local tile replay
preserve this extension and its Hermitian transpose.

The experimental Python API supports complex time histories, online spectral
observations, CPU geometry gradients and DRAM or file-backed field banks.
CUDA staging uses bounded pinned buffers. Interior tiles skip identity phase
multiplication and use read-only slices until the owned transport packet is
formed. Endpoint adjoints transfer only owned cells. Device-side zeroing
initializes halo seeds before the transpose. These changes reduce copies
without discarding any physical field update or halo derivative.

On an RTX 5880 Ada, a $256\times96\times96$ complex FP64 grid with 32 steps,
tile width 32 and temporal depth 8 gave a median complete iteration of
1.592 seconds with two asynchronous slots over three repetitions after
warmup. The resident fused path took 0.277 seconds. Peak Torch CUDA
allocation was 609,492,480 bytes for streaming and 1,080,850,944 bytes for
residency. The streamed gradient relative $L_2$ difference was
$1.42\times10^{-16}$. Thus streaming saves approximately 43.6 percent of
device allocation but costs a factor of 5.74 in runtime for this case.
The problem fits in physical VRAM and the run is short. These measurements
do not establish an external-solver advantage or beyond-VRAM capacity.
Reported transfer bytes are logical payload counters, not physical PCIe
or NVMe traffic. Mixed precision, asynchronous storage prefetch and
complete large complex-domain forward and adjoint validation remain future work.

\section{Full spectral objective and initial capacity evidence}
The supplied color-router schedule contains nine wavelengths and sixteen
pupil rays with an unnormalized total ray weight of 0.996198318172399.
At a 50 nm mesh and 1600 time steps, the relaxed seed gives an information
objective of 1.236263348891092 bits per pixel. The complete density-gradient
norm is 0.029781743905171024. Sequential case replay on an RTX 3060 reaches
a peak Torch CUDA allocation of 328,878,592 bytes. The compact reference
cache holds 42,471,936 bytes. Total forward and backward time is 5416.329
seconds, including overlapping development activity during part of the run.
This elapsed time is diagnostic rather than an isolated performance result.
The corresponding TORCWA order-eight objective is 1.2064319578769425.
The difference is a numerical discrepancy for one fixed design, not an
optimization improvement. Optical mesh, duration and Fourier-order convergence
and subsequent design optimization remain necessary.

A separate $1024\times1024\times576$ complex FP64 capacity experiment has
57,982,058,496 bytes of electric and magnetic fields alone, exceeding the
device's 51,526,500,352 bytes of physical VRAM. A completed ten-step retry used
673.955 seconds for forward execution and 3121.894 seconds through backward
and final checks. Point histories matched a smaller full-autograd domain
whose boundaries lie outside the finite dependency cone. The maximum gradient
discrepancy in that domain was $3.39\times10^{-21}$, and the finite global
gradient norm matched the local oracle at $4.0470\times10^{-5}$.
Peak Torch CUDA allocation was 7,091,686,400 bytes and sampled process RSS
was 22,506,119,168 bytes. Backward held at most 175,519,039,488 logical file
bytes and all scratch banks were released. These are short capacity and VJP
consistency results, not a converged optical application or speed advantage.
Torch allocations exclude context memory and process RSS excludes OS file
cache. The initial attempt was interrupted by a workstation restart and its
forward-only snapshot is retained separately from the completed retry.
\section{Coupled dispersive material transpose}
A separate resident operation now includes passive Drude and Lorentz poles.
The pole parameters are spatially diagonal tensors and may depend on a
Torch geometry map. Define dimensionless parameters $A_j=(\omega_{0j}\Delta t)^2$,
$G_j=\gamma_j\Delta t$ and $S_j=f_j\Delta t^2$, where $f_j$ is the oscillator
strength. With $Q_j=\Delta t\,\partial_t P_j$, define
\begin{equation}
 a_j=A_j/2,\qquad d_j=1+G_j/2+A_j/4,\qquad k_j=S_j/(4d_j).
\end{equation}
All products below act at individual Yee components. Let $c$ be the
electric curl increment before division by permittivity. The coupled
trapezoidal electric update is
\begin{align}
 r_j&=(Q_j-a_jP_j)/d_j, & b&=\epsilon_\infty+\sum_j k_j,\\
 \widetilde E&=\frac{(\epsilon_\infty-\sum_j k_j)E+c-\sum_j r_j}{b}, &
 \delta_j&=r_j+k_j(\widetilde E+E),\\
 P'_j&=P_j+\delta_j, & Q'_j&=-Q_j+2\delta_j.
\end{align}
The electric source is added after this correction. The magnetic update
uses that injected electric field. Source values are fixed in this API.
The new material states join E, H and all CPML memories in every checkpoint.

After transposing the magnetic update and observation, let $\bar e$ be the
incoming electric adjoint. For complex Bloch fields, use Hermitian field
transposes and real material derivatives. Set
\begin{align}
 t_j&=\bar P'_j+2\bar Q'_j, &
 u&=(\bar e+\sum_j k_jt_j)/b,\\
 v&=-\operatorname{Re}[(\bar e+\sum_j k_jt_j)^*\widetilde E]/b, &
 \bar r_j&=t_j-u.
\end{align}
The material and local state derivatives are
\begin{align}
 \bar\epsilon_\infty&=\operatorname{Re}(u^*E)+v,\\
 \bar k_j&=\operatorname{Re}[t_j^*(\widetilde E+E)-u^*E]+v,\\
 \bar a_j&=-\operatorname{Re}(\bar r_j^*P_j)/d_j,\\
 \bar d_j&=-\operatorname{Re}(\bar r_j^*r_j)/d_j-\bar k_j k_j/d_j,\\
 \bar S_j&=\bar k_j/(4d_j),\qquad
 \bar A_j=\bar a_j/2+\bar d_j/4,\qquad \bar G_j=\bar d_j/2,\\
 \bar P_j&=\bar P'_j-a_j\bar r_j/d_j,\qquad
 \bar Q_j=-\bar Q'_j+\bar r_j/d_j,\\
 \bar E&=\sum_j k_jt_j+(\epsilon_\infty-\sum_j k_j)u.
\end{align}
The curl adjoint receives $u$ with the existing Courant factor. Torch
differentiates the timestep scaling and preceding geometry map. This
construction retains a bounded set of physical states rather than a
computational graph across all time steps.

Shared material parameters retain their original tensor shapes in a compact
input buffer. Broadcasting supplies the field update without storing dense
copies of a uniform pole rate. The transpose reduces cell and component
derivatives to those original shapes before temporal accumulation. A scalar
shared by all poles also receives the pole-axis sum. This changes floating
point summation order but not the represented material or storage precision.
The spatial P and Q states remain necessary, and this parameter reduction
does not imply sparse field evolution or a bound on the full solver peak.

The reference implementation uses Torch operations and an explicit analytic
transpose. Small CPU tests cover FP32 and FP64, fixed Bloch phase, scalar
and diagonal epsilon, all four material parameter derivatives, checkpoint
restoration and online spectral observations. Native Drude, Lorentz and
multipole forward traces agree with the new recurrence. A full-time
autograd oracle and directional tests check the material VJP. Fixed
collocated planes connect it to reference-normalized flux objectives.
The material suite also passes on an RTX 3060, including FP32 and FP64
VJPs and asynchronous device, host and file checkpoint restoration.
Sparse oscillator state and physical gradient convergence remain pending.
An experimental spatial extension is described below.
These tests do not establish a new dispersive integration algorithm or
a large-domain performance advantage.

\subsection{Native material transpose and shared-parameter reductions}
The resident CUDA path fuses the electric curl and coupled material update.
Each component first sums the pole responses, solves for
$\widetilde E$, and updates P and Q in place. The magnetic kernel follows
electric source injection. Backward recomputes $\widetilde E$ without
changing the primal restart state. Magnetic and electric curl transposes
surround a pointwise material transpose. This keeps the existing Yee and
CPML derivative conventions and avoids a time-axis computational graph.

Spatial material derivatives have a unique writer for each cell or
component. Shared pole rates instead use a fixed-order reduction within
256-thread blocks. Block partials accumulate across reverse time steps,
followed by a compact reduction at the end of backward. For $m$ shared
parameters and $N$ cells, this temporary contains
\begin{equation}
 m\left\lceil\frac{N}{256}\right\rceil
\end{equation}
real values rather than a dense pole-by-cell material-gradient array.
This reduction uses no atomic updates. The P/Q states and their adjoints
remain dense, and this count does not bound total solver memory.

The native path is tested against the full-time Torch graph for real and
complex fields, scalar and diagonal epsilon, spatial and shared material
parameters, and FP32 and FP64 arithmetic. A complete one-step state
transpose test includes nonzero material and CPML endpoint adjoints.
Additional checks cover fixed spectral planes, repeated backward,
mixed-tier checkpoint restoration and a 64-pole Drude limit.
This is an implementation of the stated discrete transpose, not a claim
of a new dispersive time integrator.

On an RTX 5880 Ada, a matched resident comparison uses two poles,
three device checkpoints, one warmup per execution mode and three timed
repetitions with rotating order. Every repetition checks signals and material
gradients against the Torch implementation. At $128^3$ cells and 128 steps,
the median results are shown in Table~\ref{tab:ade-cuda}.
Full wall includes parameter packing, forward, the objective, checkpoint
replay, backward and shared-gradient reduction. Model construction, input
allocation and final host copies are excluded. Peak memory counts Torch
allocations and excludes the CUDA context. These measurements do not establish
external-solver performance, optical convergence or out-of-core ADE capacity.
\begin{table}[ht]
\centering
\begin{tabular}{lrrrr}
\toprule
Fields & Torch (s) & Fused (s) & Ratio & Peak GiB, Torch / fused \\
\midrule
Real FP32 & 4.2505 & 0.5307 & 8.01 & 1.572 / 0.804 \\
Complex FP64 & 19.9781 & 4.6447 & 4.30 & 5.512 / 3.129 \\
\bottomrule
\end{tabular}
\caption{Native resident ADE forward and backward ablation at $128^3$ cells.
Fused backward alone gives ratios of 1.28 in both cases because its replay
retains Torch forward. Raw repetition records and source hashes accompany
the repository validation report.}
\label{tab:ade-cuda}
\end{table}
\subsection{Material states in space-time slabs}
The spatial extension retains the same pointwise ADE recurrence.
Consequently, it adds no spatial dependency beyond the
$K$ halo for $K$ complete Yee steps. Every tile reads an immutable
block-start state, evolves its extended domain, and commits only owned cells.
E/H, CPML and P/Q participate in both block and local checkpoint replay.

Global P/Q arrays have layout $(N_x,P,N_y,N_z,3)$ for contiguous x-slab
file access. Packing supplies the pole-first layout to the native CUDA
kernel. Bloch extension multiplies all replicated field and polarization
states by the winding phase. Its transpose applies the conjugate phase
before reducing overlapping halo contributions. Real material coefficients
receive no Bloch phase. Their spatial derivatives sum over replicated x
indices, while shared pole or scalar derivatives sum over all tiles.

The initial global P/Q state uses scalar-backed zero views. Evolving
states use DRAM or explicitly budgeted transient files, and CUDA scratch
is reused per staging slot. Admission includes pole states, material-gradient
carriers, local replay states and transfer buffers. The pole-dependent
workspace estimate is conservative and still requires large-run calibration.
The spectral observer and fixed-plane flux interface share this execution
path without retaining the complete signal history.

Complete block-Jacobian tests compare against resident full-time autograd,
including nonzero polarization and CPML states and arbitrary endpoint
adjoints. Tests also cover repeated Bloch windings, shared and spatial
material parameters, file-backed replay and asynchronous staging. These
checks establish discrete consistency for the tested cases. They do not
establish optical convergence, large dispersive capacity or a speed advantage.
\subsection{Bounded material-aware policy calibration}
Policy calibration now covers both dielectric and coupled ADE execution,
using either point histories or online point spectra. Generated candidates
vary spatial width, temporal depth and checkpoint capacity. A budget check
includes the complete requested duration and the tuner comparison workspace.
If a generated tile does not fit, a monotone search reduces its core width
at fixed temporal depth. Depth is reduced only when a one-cell core is still
inadmissible. File fallback requires an explicitly supplied directory and
capacity. Explicit candidate lists are not modified.

Calibration previously retained a complete reference gradient for each
duration. The current implementation uses a byte-bounded reference cache.
Oversized references bypass retention and are recomputed for comparison.
Compact owned copies prevent a gradient view from retaining a larger packed
material allocation. The reservation also includes the simultaneous candidate,
reference and conversion workspace. Full host material-gradient buffers remain
necessary, even with zero cache capacity. These checks do not include
caller-owned geometry graphs or optimizer storage.

For meaningful material checks, derivatives are compared in the coordinates
\begin{equation}
 \vartheta=(\epsilon_\infty,s\Delta t^2,\omega_0\Delta t,\gamma\Delta t).
\end{equation}
For example, the strength derivative is divided by $\Delta t^2$ before
comparison. An absolute tolerance applied directly to an SI strength
derivative could otherwise accept a substantial relative error.
The Drude limit requires no division by $\omega_0$.
Spectral calibration divides the DFT fields by the prefix duration before
forming the proxy loss. Frequencies remain fixed and a supplied window is
sliced to each prefix.

Each measured candidate uses the same two-duration replay cost model, with
optional longer calibration. Timing includes material packing, solver
preparation, forward, the proxy objective, replay and backward. Gradient
comparison, model construction, the caller's geometry graph and optimizer
are excluded. Tests verify derivative disagreement even when forward values
are unchanged, bounded cache ownership and admission before packing.
This remains a bounded empirical search. Full-duration ranking on large
dispersive applications, physical storage bandwidth and resident-versus-streamed
selection require separate experiments.
\end{document}
